\documentclass[11pt]{article}
\usepackage[a4paper,margin=26mm]{geometry}
\usepackage{amsmath,amssymb,amsthm,hyperref}
\newtheorem{theorem}{Theorem}
\newtheorem{lemma}{Lemma}
\newtheorem{proposition}{Proposition}
\newcommand{\E}{\mathbb E}
\newcommand{\Pol}{\operatorname{Pol}}
\title{A localized Laguerre bootstrap toward a quadratic per-die bound}
\author{Research proof: three independent mathematical reviews completed}
\date{30 September 2026}
\begin{document}\maketitle
\noindent\textbf{Status.} Three independent full mathematical reviews
completed, with no gap identified; this is not formal verification.
The classical Laguerre input was checked independently in the primary
paper. The proof uses the reviewed comparison machinery of
\path{../power_per_die_lower_bound.tex}.
No assertion of literature novelty is made.

\begin{theorem}\label{main}
For every fixed $\varepsilon>0$ there is $c_\varepsilon>0$ such that
every individual die in a tie-free permutation-fair family of $n$
equiprobable finite dice has at least $c_\varepsilon n^{2-\varepsilon}$
faces. More generally, every monotone mixed-moment tensor with $m$
columns has first and last row weights at most
$C_\varepsilon m^{-2+\varepsilon}$.
\end{theorem}

Throughout, assume $p_j(q)\in[0,1]$ are nondecreasing in $q$, row weights
$w_q>0$ sum to one, and
$\E_w\prod_{j\in S}p_j=1/(|S|+1)$ for every $S\subseteq[m]$.
Use the notation of the cited proof:
\[
 u_j=1-p_j,\quad s=\sum_j u_j,\quad
 \ell=\lfloor m/4\rfloor,\quad N=m-\ell,\quad
 \alpha=\ell/m,\quad x=\alpha s,\quad\theta=1/16.
\]
For a uniform $\ell$-subset $B$, let
$A_B=\prod_{j\in B}p_j$, $\bar A=\E_BA_B$, and
$v_B=N^{-1}\sum_{j\notin B}\ell u_j$.
The probability measure $\mu$ on $x$ has row masses
$(\ell+1)w_q\bar A(q)$, and $\beta$ has density
\[
 b_\ell(z)=\frac{\ell+1}{\ell}(1-z/\ell)^\ell
                 \boldsymbol1_{[0,\ell]}(z).
\]
The finite positive measure $\nu$ on $x$ has row masses $mw_q$;
its total mass is $m$. Constants may depend on explicitly fixed
exponents or derivative orders, but never on $m$ or the tensor.

\section*{Inputs already proved in the companion notes}
We use the following facts, all proved in the cited power-bound note
and \path{../rational_designs/asymptotic/weighted_success_subset_moments.tex}:
\begin{align}
 \max_j u_j&\le C(s+1)/m,&
 m\Pr_w(s\le t)&\le C(t+1),\label{inputs}\\
 \bar A&\le e^{-\alpha s},&
 \E_w(s+1)^k e^{-a s}&\le C_a^{k+1}k!/m.\label{inputs2}
\end{align}
For $s\le m/256$, all $p_j\ge1/2$. Under the success-weighted
subset law $\Pr(B)=A_B/e_\ell(p)$, $\Delta=v_B-x$ satisfies
\begin{equation}\label{weighted}
 0\le\E\Delta\le C s(s+1)/m,\qquad
 \E\Delta^2\le C\{s(s+1)/m+s^2(s+1)^2/m^2\}.
\end{equation}
For any polynomial of degree at most $N$,
\begin{equation}\label{exact}
 \beta(P)=(\ell+1)\E_w\E_B
           [A_B\Pol_N(P)((\ell u_j)_{j\notin B})].
\end{equation}
If $P=Q_1Q_2$ is a product of degree-$d$ Laguerre sums with coefficient
$\ell^1$ norms $B_1,B_2$, the Gaussian derivative and centered-variance
bounds give the following useful tail estimate. The part of the exact
polarization Taylor series of orders $k\ge J$ has integrated absolute
value at most
\begin{equation}\label{gaussiantail}
 C B_1B_2(Cd/m)^{J/2},
\end{equation}
provided $d=o(m)$, $2d\le N$, and $J\ge2$ is fixed and even.
Indeed the proof in the companion note gives, before summing,
\[
 C B_1B_2\sum_{k=J}^{2d}(C\sqrt{d/m})^k
       k^{k/2}\lceil k/2\rceil!/k!,
\]
and the factorial ratio is bounded by $C^k(k+1)$.
Its geometric sum proves \eqref{gaussiantail}; the earlier restriction
$d\le\sqrt m$ is unnecessary for this tail estimate.

\section*{The classical localized Laguerre estimate}
For fixed $a\ge0$, the standard normalized Laguerre functions
\[
 \mathcal L_r^a(y)=\sqrt{\frac{r!}{\Gamma(r+a+1)}}
                    L_r^{[a]}(y)y^{a/2}e^{-y/2}
\]
satisfy
\[
 |\mathcal L_r^a(y)|\le C_a
 \begin{cases}
 (y\rho)^{a/2},&0<y\le1/\rho,\\
 (y\rho)^{-1/4},&1/\rho<y\le\rho/2,
 \end{cases}
 \qquad\rho=4r+2a+2.
\tag{L}
\]
These are the standard bounds stated explicitly in Plewa
\cite[Section 2, immediately before (1)]{Plewa}, with sources
Muckenhoupt \cite[p.~435]{Muckenhoupt} and Askey--Wainger
\cite[p.~699]{AskeyWainger}. Only fixed nonnegative integer parameters
$a$ are used here. Equivalently, uniformly for $0\le y\le\rho/2$,
\begin{equation}\label{lagloc}
 |L_r^{[a]}(y)|\le C_a e^{y/2}(r+1)^a
                       (1+(r+1)y)^{-a/2-1/4}.
\end{equation}

Put
\[
 T_d(y)=\sum_{j<d}(1-j/d)L_j(y)=\frac1dL_{d-1}^{[2]}(y),
 \quad U_d(z)=d^{-2}T_d(\theta z)^2,
 \quad R_d(z)=(2/(\theta d)-z)T_d(\theta z)^2.
\]
For every fixed integer $J$, all $0\le k\le J$, all sufficiently
large $d$, and $0\le z\le2d$, \eqref{lagloc} and
$D^hT_d=(-1)^hL_{d-h-1}^{[h+2]}/d$ give
\begin{align}
 |U_d^{(k)}(z)|&\le C_J e^{\theta z}d^k
                      (1+dz)^{-k/2-5/2},\label{Uder}\\
 |R_d^{(k)}(z)|&\le C_J e^{\theta z}d^{k+1}
                      (1+dz)^{-k/2-3/2}.\label{Rder}
\end{align}
Here $d-h-1$ is comparable to $d$, and $\theta z\le d/8<\rho/2$.
The product rule proves \eqref{Uder}; multiplication by
$2/(\theta d)-z$ proves \eqref{Rder} using
$2/(\theta d)+z\le C(1+dz)/d$.
The global Laguerre coefficient products $B_1B_2$ are $O(1)$ for
$U_d$ and $O(d)$ for $R_d$, by the exact identity
$yT_d=d^{-1}\sum_{j<d}L_j-L_d$.

\section*{A lower-tail estimate for subset averages}
\begin{lemma}\label{tail}
On every row, for a uniform subset $B$,
\[
 \Pr_B(v_B<x/2)\le\exp[-c m\min(x,1)].
\]
This statement concerns the uniform subset law. In expressions with
success weighting, the unnormalized bad contribution obeys the same
bound because $0\le A_B\le1$.
\end{lemma}
\begin{proof}
The whole population $z_j=\ell u_j$ has mean $x$ and is contained in
$[0,R]$ with $R=C(s+1)$ by \eqref{inputs}. The complement is a uniform
$N$-subset. Maclaurin's inequality and the chord bound for $e^{-z/R}$
give
\[
 \E_B e^{-Nv_B/R}
 \le\left(\frac1m\sum_j e^{-z_j/R}\right)^N
 \le\exp[-Nx(1-e^{-1})/R].
\]
Markov's inequality at $v_B=x/2$ now yields
$\exp[-(1-e^{-1}-1/2)Nx/R]$.
Since $N\asymp m$ and $s=x/\alpha\asymp x$, this is the claimed bound.
The row $x=0$ causes no difficulty.
\end{proof}

\section*{Localized comparison with a distribution-function hypothesis}
\begin{proposition}\label{localized}
Fix $0<\eta<b<1$ and a positive constant $C_0$. Suppose
\begin{equation}\label{cdf}
 \nu([0,t])\le C_0(t+h)\qquad(t\ge0),\quad0<h\le1.
\end{equation}
For $m^\eta\le d\le m^b$, a fixed even $J\ge4$, and any fixed
$A>0$, one has, uniformly in this range,
\begin{align}
 |\mu(U_d)-\beta(U_d)|
 &\le C_J\{m^{-1}(1+hd)+(Cd/m)^{J/2}\}+O(m^{-A}),\label{Ucompare}\\
 |\mu(R_d)-\beta(R_d)|
 &\le C_J\{(d/m)(1+hd)+d(Cd/m)^{J/2}\}+O(m^{-A}).\label{Rcompare}
\end{align}
The implied constants may depend on $\eta,b,A,C_0,J$.
\end{proposition}
\begin{proof}
Write $\tau=0$ for $U_d$, $\tau=1$ for $R_d$, and let
$b_0=5/2$ or $3/2$, respectively. Thus the coefficient product is
$O(d^\tau)$ and the fixed-order derivative bounds have the form
\[
 |P^{(k)}(z)|\le C_J e^{\theta z}d^{k+\tau}
                         (1+dz)^{-k/2-b_0}.
\tag{D}
\]
We compare exact polarization to $P(v_B)$, and then $P(v_B)$ to $P(x)$.

First consider $0\le x\le1/d$. The global Gaussian-variance argument
in the companion proof bounds the full polarization error pointwise
by $C d^{\tau+1}/m$, after averaging with $A_B$; here $s+1=O(1)$
and its geometric series starts at order two. Its sharp weighted
Taylor argument bounds the averaging error by the same quantity,
using \eqref{weighted}; the factors singular in $x$ are canceled by
the factors $s$ in those moment bounds. At $x=0$ both errors vanish.
Consequently the integrated contribution is at most
\[
 C\frac{d^{\tau+1}}m\nu([0,1/d])
 \le C\frac{d^\tau}m(1+hd).
\tag{S}
\]

Next consider $1/d<x\le d$. This region lies in $s\le m/256$ for
large $m$ because $d\le m^b=o(m)$. Call a subset good if $v_B\ge x/2$.
Always $v_B\le x/(1-\alpha)<2x$. On good subsets the centered variance
satisfies
\[
 \sum_{j\notin B}(z_j-v_B)^2\le C N(s+1)v_B.
\]
Combining (D) at $v_B$ with the centered elementary-coefficient bound,
each Taylor term $2\le k<J$ is at most
\[
 C_J d^{\tau+k/2}m^{-k/2}(s+1)^{k/2}
                 e^{\theta v_B}(1+dx)^{-b_0}.
\]
After averaging with $A_B$, its exponential factor is at most
$e^{-c x}$; the fixed power of $s+1$ can be absorbed by reducing $c$.
Since $d/m\le1$, summing these fixed orders bounds the result by
\begin{equation}\label{envelope}
 C_J\frac{d^{\tau+1}}m e^{-c x}(1+dx)^{-b_0}.
\end{equation}
The contribution from orders $k\ge J$, integrated over all rows,
is bounded separately by \eqref{gaussiantail}.

On bad subsets, the global Gaussian-variance bound for each of these
finitely many orders is a polynomial in $m$ times $e^{Cx}$.
Lemma~\ref{tail}, with $A_B\le1$, bounds its average by a polynomial
in $m$ times
$e^{Cx-cm\min(x,1)}$. For $1/d<x\le d$, this is
$O(m^{-A})$ for every fixed $A$, uniformly: when $x\le1$ use
$m/d\ge m^{1-b}$, and when $x\ge1$ use $d=o(m)$.
The additional integration against a measure of total mass $m$
does not change this conclusion.

For the weighted averaging error write the exact Taylor remainder
at $x$ and retain the full mean term $P'(x)\E\Delta$.
On good subsets the intervening points lie in $[x/2,2x]$.
Using (D) with $k=1,2$ and \eqref{weighted}, then multiplying by
$\bar A$, bounds this contribution by \eqref{envelope}.
Explicitly, before the exponential factor is absorbed, the mean term
has size at most
\[
 C_J d^{\tau+1}(1+dx)^{-b_0-1/2}\,x(x+1)/m,
\]
and the quadratic remainder has size at most
\[
 C_J d^{\tau+2}(1+dx)^{-b_0-1}
       \{x(x+1)/m+x^2(x+1)^2/m^2\}.
\]
The former is bounded as claimed because
$x/\sqrt{1+dx}\le1$ for $x\le d$; for the latter use
$dx/(1+dx)\le1$. Fixed remaining powers of $x+1$ are absorbed by
the exponential margin. On bad subsets bound the remainder by
$|P(v_B)|+|P(x)|+|P'(x)||\Delta|$, a polynomial in $m$ times $e^{Cx}$.
Multiplication by $\bar A$ converts its weighted probability exactly
back into $\E_B A_B\boldsymbol1_{\rm bad}$, bounded by the uniform
bad probability. Thus Lemma~\ref{tail} again gives $O(m^{-A})$.

For $x>d$ the global estimates suffice. In the integrated polarization
series reserve half of the negative exponential factor to obtain
$e^{-cd}$, and use \eqref{inputs2} with the other half. The same
geometric-series calculation gives a polynomial factor times $e^{-cd}$.
For averaging, use the nonsingular bound
$|P'|\le C d^{\tau+1}e^{\theta z}$ and the uniform-subset
Cauchy--Schwarz estimate
$\E_B A_B|v_B-x|\le C(s+1)e^{-\alpha s/2}/\sqrt m$.
The exponential margin and \eqref{inputs2} again give a polynomial
factor times $e^{-cd}$. Since $d\ge m^\eta$, both tails are
$O(m^{-A})$.

Finally, for $g(x)=e^{-cx}(1+dx)^{-b_0}$, Stieltjes integration by
parts and \eqref{cdf} give
\[
 \int g\,d\nu
 \le C_0\int_0^\infty(t+h)(-g'(t))\,dt
 =C_0\left(\int_0^\infty g(t)\,dt+h\right)
 \le C(1/d+h),
\]
including a possible atom at zero. Integrate \eqref{envelope}, combine
with (S), the high-order tail and \eqref{exact}, to obtain the result.
\end{proof}

\section*{The bootstrap}
\begin{lemma}\label{base}
One has $\nu([0,t])\le C(t+m^{-1/2})$ for all $t\ge0$.
\end{lemma}
\begin{proof}
The reviewed bilinear comparison gives
$\mu(U_r)\le\beta(U_r)+Cr/m\le C/r+Cr/m$ for $r\le\sqrt m$.
Here $\beta(U_r)\le C/r$ follows by Laguerre orthogonality and
$b_\ell(z)\le2e^{-\theta z}$. Also $U_r(z)\ge c$ for
$0\le z\le c_1/r$, by the elementary near-zero Laguerre estimate.
Choose $c_1$ small enough that $x\le c_1$ implies $s\le1/2$;
then $\bar A\ge1-s\ge1/2$, so $\nu$ and $\mu$ are comparable
on this interval. Taking $r\asymp\min(\sqrt m,1/t)$ proves the
claim for $t\le c_1$. For larger $t$, the original mass bound
in \eqref{inputs} is $O(t+1)=O(t)$.
\end{proof}

\begin{proposition}\label{step}
If, for a fixed $0<a<1$, every such tensor satisfies
$\nu([0,t])\le C_a(t+m^{-a})$, then with $b=(1+a)/2$ one also has
$\nu([0,t])\le C_b(t+m^{-b})$.
\end{proposition}
\begin{proof}
Use Proposition~\ref{localized} with $h=m^{-a}$ and an even $J$ so
large that $J(1-b)/2>b$. Uniformly for degrees
$m^{a/2}\le r\le m^b$, it gives
\[
 \mu(U_r)\le C/r+C\{m^{-1}+m^{-a}r/m+(Cr/m)^{J/2}\}
                 +O(m^{-A})\le C/r,
\]
where $A>b$ and $r^2\le m^{1+a}$ were used.
The same near-zero positivity and $\mu$--$\nu$ comparison as in
Lemma~\ref{base} imply $\nu([0,c_1/r])\le C/r$.
For $t\lesssim m^{-a}$ choose
$r\asymp\min(m^b,1/t)$; these degrees are at least $m^{a/2}$
for large $m$. For $t\gtrsim m^{-a}$ use the old estimate.
This proves the new distribution-function bound.
\end{proof}

Starting from $a_0=1/2$, iteration gives
$a_j=1-2^{-j-1}$. Fix one such $a$, set $b=(1+a)/2$, and now choose
$d=\lfloor\kappa m^b\rfloor$ in \eqref{Rcompare}, with $\kappa>0$
sufficiently small. Choose $J$ as in Proposition~\ref{step}.
The high-order tail is $o(1)$ and
\[
 (d/m)(1+m^{-a}d)=o(1)+O(\kappa^2).
\]
The scalar calculation from the reviewed power-bound proof gives
$\beta(R_d)\ge1/(96\theta)>0$. Hence $\mu(R_d)>0$ for sufficiently
small $\kappa$ and large $m$. Some row has $x<2/(\theta d)$, and
the last row minimizes $x$, so $x_K<2/(\theta d)$.
The already proved distribution-function estimate with exponent $b$
then gives
\[
 mw_K\le\nu([0,2/(\theta d)])\le C_b/d,
 \qquad w_K\le C_b m^{-1-b}.
\]
Reverse the rows and replace every $p_j$ by $1-p_j$ for the first row.
Since the exponents $b$ approach one, this proves Theorem~\ref{main}.
The constants are not uniform as $\varepsilon\downarrow0$; a quadratic
$\Omega(n^2)$ lower bound is not asserted.

\begin{thebibliography}{9}
\bibitem{Plewa} P. Plewa, Hardy's inequality for Laguerre expansions of
Hermite type, \emph{J. Fourier Anal. Appl.} 25 (2019), 1855--1873.
\url{https://doi.org/10.1007/s00041-018-9642-2}.
\bibitem{Muckenhoupt} B. Muckenhoupt, Mean convergence of Hermite and
Laguerre series. II, \emph{Trans. Amer. Math. Soc.} 147 (1970), 433--460.
\bibitem{AskeyWainger} R. Askey and S. Wainger, Mean convergence of
expansions in Laguerre and Hermite series, \emph{Amer. J. Math.} 87
(1965), 695--708.
\end{thebibliography}
\end{document}
